\documentclass[review,3p]{elsarticle}

%% ============================================================
%% PAPER 1: TOPOLOGY ABSTRACTION EFFECTS
%% Target Journal: Energy Conversion and Management
%% Version: 3.0 — Bereinigt, Appendices korrigiert, Ergebnisse ausstehend
%% ============================================================

\usepackage[utf8]{inputenc}
\usepackage[T1]{fontenc}
\usepackage{amsmath,amssymb,amsfonts}
\usepackage{graphicx}
\usepackage{booktabs}
\usepackage{multirow}
\usepackage{array}
\usepackage{siunitx}
\usepackage{xcolor}
\usepackage{hyperref}
\usepackage{cleveref}
\usepackage{lineno}
\usepackage{float}
\usepackage{subcaption}
\usepackage{enumitem}
\usepackage[version=4]{mhchem}

\modulolinenumbers[5]
\sisetup{per-mode=symbol, bracket-unit-denominator=false}
\hypersetup{colorlinks=true, linkcolor=blue, citecolor=blue, urlcolor=blue}

\newcommand{\placeholder}[1]{\textcolor{red}{[#1]}}

\journal{Energy Conversion and Management}

\begin{document}

\begin{frontmatter}

%% ============================================================
%% TITEL
%% ============================================================
\title{Topology Abstraction Effects on Dispatch Optimization of
Electrified District Heating Networks: A Controlled MILP Comparison}

\author[tum]{Lukas [Nachname]\corref{cor1}}
\ead{lukas.nachname@tum.de}
\author[tum]{[Betreuer]}
\author[partner]{[Industriepartner]}

\cortext[cor1]{Corresponding author.}

\affiliation[tum]{organization={Lehrstuhl für Energiewirtschaft und Anwendungstechnik,
                Technische Universität München},
                city={München},
                country={Germany}}

\affiliation[partner]{organization={[Industriepartner Institution]},
                city={[Stadt]},
                country={Austria}}

%% ============================================================
%% ABSTRACT
%% ============================================================
\begin{abstract}
District heating network planning studies require practitioners to choose
a topology representation level, yet no quantitative guidance exists for
this decision. Prior comparisons confound topology abstraction with
physics model changes, rendering it impossible to attribute observed
cost differences to either factor independently.

This paper presents the first controlled comparison that isolates
topology as the sole experimental variable. Three abstraction levels
are defined within a unified mixed-integer linear programming (MILP)
framework: L1 (copperplate, no network physics), L2 (simplified
multi-node with physics-based thermal losses), and L3 (detailed
multi-node with identical physics to L2). The framework is validated
against measured operational data from an Austrian industrial heat
network comprising heat pumps, thermal storage, a gas boiler, an
electrode boiler, and grid electricity coupling. Generalizability is
assessed across 36 synthetic network configurations spanning systematic
variations in pipe length, demand heterogeneity, node count, and
storage capacity.

Results quantify the topology effect: the L1$\rightarrow$L2 transition
(physics inclusion) drives cost differences of \placeholder{X}\%
(\placeholder{€Y,000}/year) and \ce{CO2} differences of
\placeholder{Z}~t/year, while the L2$\rightarrow$L3 transition
(spatial refinement) adds \placeholder{$<$X}\% for uniform-demand
networks but up to \placeholder{Y}\% for spatially heterogeneous
systems. Thermal storage dispatch emerges as the primary diagnostic
variable for model adequacy: L1 systematically misrepresents charging
patterns by ignoring loss-dependent temporal holding costs, shifting
optimal charge timing by \placeholder{X}~hours on average in winter.
The synthetic analysis reveals that the L1$\rightarrow$L2 gap scales
with total pipe length ($R^2\!=\,$\placeholder{XX}), while the
L2$\rightarrow$L3 gap is driven by demand heterogeneity. These
findings yield actionable selection criteria: L2 is the minimum
fidelity for storage-containing systems, while L3 is warranted when
the demand heterogeneity index exceeds \placeholder{0.X}.
\end{abstract}

\begin{keyword}
District heating \sep Mixed-integer linear programming \sep
Topology abstraction \sep Thermal storage dispatch \sep
Heat pump optimization \sep Network modeling fidelity
\end{keyword}

\end{frontmatter}

\linenumbers

%% ============================================================
%% NOMENCLATURE
%% ============================================================
\section*{Nomenclature}

\subsection*{Sets and Indices}
\begin{tabular}{@{}p{2.8cm}p{9.5cm}@{}}
$t \in \mathcal{T}$               & Time steps, $t = 1, \ldots, |\mathcal{T}|$ \\
$n \in \mathcal{N}$               & Network nodes \\
$\mathcal{N}^{\mathrm{prod}}$     & Producer nodes (generation/storage) \\
$\mathcal{N}^{\mathrm{con}}$      & Consumer nodes (demand only) \\
$p \in \mathcal{P}$               & Pipe segments \\
$g \in \mathcal{G}$               & Thermal generators (all types) \\
$h \in \mathcal{H}$               & Heat pumps \\
$s \in \mathcal{S}$               & Thermal storage units \\
$r \in \mathcal{R}$               & Power-to-heat units \\
$\mathcal{G}_n$                   & Generators at node $n$ \\
$\mathcal{S}_n$                   & Storage units at node $n$ \\
\end{tabular}

\subsection*{Parameters}
\begin{tabular}{@{}p{2.8cm}p{2.2cm}p{7.5cm}@{}}
$\Delta t$                              & h               & Time step duration (1~h unless stated) \\
$U_p$                                   & W/(m$\cdot$K)   & Linear heat transfer coefficient, pipe $p$ \\
$L_p$                                   & m               & Length of pipe $p$ \\
$T_{\mathrm{ground}}(t)$               & \si{\celsius}   & Ground temperature at time $t$ \\
$T_{\mathrm{sup}}^{\mathrm{nom}}(t)$   & \si{\celsius}   & Nominal supply temperature (heating curve) \\
$T_{\mathrm{ret}}^{\mathrm{nom}}$      & \si{\celsius}   & Nominal return temperature \\
$\mathrm{COP}_{h,t}^{\mathrm{wrg}}$   & --              & Pre-computed COP, HP $h$, WRG source \\
$\mathrm{COP}_{h}^{\mathrm{def}}$      & --              & Default-source COP, HP $h$ (constant) \\
$\eta_g^{\mathrm{th}}$                 & --              & Thermal efficiency of generator $g$ \\
$\eta_g^{\mathrm{el}}$                 & --              & Electrical efficiency of CHP unit $g$ \\
$\eta_{s}^{\mathrm{ch/dis}}$           & --              & Storage charge/discharge efficiency \\
$\alpha_s$                             & h$^{-1}$        & Storage self-discharge rate \\
$a_{s,t}$                              & $\{0,1\}$       & Storage availability parameter \\
$\lambda_t^{\mathrm{buy/sell}}$        & \euro/MWh       & Electricity buy/sell price at time $t$ \\
$\lambda_g^{\mathrm{fuel}}$            & \euro/MWh       & Fuel price of generator $g$ \\
$\lambda^{\ce{CO2}}$                   & \euro/t\ce{CO2} & Carbon price \\
$\lambda^{\mathrm{dem}}$               & \euro/MW/a      & Annual demand charge rate \\
$e_t^{\mathrm{grid}}$                  & kg/MWh          & Hourly grid emission factor \\
$e_g^{\mathrm{fuel}}$                  & kg/MWh          & Fuel emission factor \\
$M^{\mathrm{grid}}$                    & MW              & Big-M constant for grid buy/sell logic \\
$\bar{Q}_g$                            & MW              & Installed thermal capacity of generator $g$ \\
$\bar{E}_s$                            & MWh             & Energy capacity of storage $s$ \\
$\bar{P}_s$                            & MW              & Power capacity of storage $s$ \\
$\phi_h^{\mathrm{min}}$                & --              & Minimum part-load ratio of heat pump $h$ \\
$f^{\mathrm{yr}}$                      & --              & Year fraction: $|\mathcal{T}|\Delta t / 8760$ \\
$\mathrm{HI}$                          & --              & Demand heterogeneity index (Eq.~\ref{eq:hi}) \\
\end{tabular}

\subsection*{Decision Variables}
\begin{tabular}{@{}p{2.8cm}p{2.2cm}p{7.5cm}@{}}
$P_t^{\mathrm{buy/sell}}$         & MW        & Grid electricity import/export \\
$P^{\mathrm{peak}}$               & MW        & Peak grid import (annual maximum) \\
$Q_{g,t}^{\mathrm{th}}$          & MW        & Thermal output of generator $g$ \\
$P_{g,t}^{\mathrm{el}}$          & MW        & Electrical output of CHP unit $g$ \\
$P_{h,t}^{\mathrm{el}}$          & MW        & Electrical input to heat pump $h$ \\
$Q_{h,t}^{\mathrm{wrg/def}}$     & MW        & HP $h$ output from WRG/default source \\
$Q_{s,t}^{\mathrm{ch/dis}}$      & MW        & Storage $s$ charge/discharge power \\
$E_{s,t}$                        & MWh       & Storage $s$ state of charge \\
$F_{g,t}$                        & MW        & Fuel input of generator $g$ \\
$Q_{n,t}^{\mathrm{dump}}$        & MW        & Heat curtailment at node $n$ \\
$Q_{p,t}$                        & MW        & Heat flow in pipe $p$ \\
$u_{h,t}$                        & $\{0,1\}$ & Binary: heat pump $h$ online/offline \\
$z_t$                            & $\{0,1\}$ & Binary: grid import mode ($z=1$) \\
$\delta_{s,t}^{\mathrm{ch/dis}}$ & $\{0,1\}$ & Binary: storage charge/discharge mode \\
\end{tabular}

%% ============================================================
%% 1. INTRODUCTION
%% ============================================================
\section{Introduction}
\label{sec:introduction}

\subsection{Motivation}
\label{sec:motivation}

The decarbonization of heating and cooling represents one of the
largest unresolved challenges in the European energy transition.
Space and process heat account for approximately 50\% of final
energy consumption in the European Union, of which roughly
60\% is still provided by fossil fuels~\cite{EUheating2016}.
District heating (DH) networks offer a systemic pathway for
decarbonization through sector coupling: large-scale heat pumps
converting waste heat and ambient sources, combined heat and
power (CHP) units providing flexibility, thermal energy storage
enabling temporal load shifting, and electrode boilers absorbing
surplus renewable electricity~\cite{Lund2014_4GDH,Werner2017}.
The fourth-generation paradigm (4GDH) envisions networks operating
at low supply temperatures to maximize heat pump efficiency and
integrate low-grade renewable sources~\cite{Lund2014_4GDH,Buffa2019}.

Planning and operational optimization of such multi-source networks
increasingly relies on mixed-integer linear programming (MILP)
due to its global optimality guarantees, mature solver technology,
and tractable handling of binary commitment decisions. However,
every MILP-based planning study confronts a fundamental decision
that has received remarkably little systematic attention: the level
of spatial detail at which the network is represented.

This choice spans a wide spectrum. At one extreme,
\emph{copperplate} models aggregate all components and demands
to a single bus, ignoring spatial structure entirely. At the
other extreme, \emph{full-topology} models resolve the physical
pipe-by-pipe network, capturing thermal losses that typically
constitute 5--15\% of annual heat supply in Central European
networks~\cite{Frederiksen2013,Talebi2016}.
Between these poles lies a continuous range of aggregation levels,
each representing a different trade-off between physical fidelity
and computational tractability.

The practical consequences of this choice affect three
interlinked dimensions. First, cost estimation accuracy: simplified
models may systematically underestimate operational costs by
ignoring the thermal generation overhead required to compensate
pipe losses. Second, dispatch quality: optimal scheduling decisions
for heat pumps, storage, and boilers depend on energy balances that
are distorted when losses are absent or misrepresented. Third,
\ce{CO2} accounting: because network losses require additional
generation from sources with non-zero emission factors,
topology-induced generation biases translate directly into
emission accounting errors.

Despite the centrality of this decision, practitioners currently
lack quantitative guidance on the accuracy--tractability trade-off
across topology levels. The present paper addresses this gap
through a controlled experimental comparison.

\subsection{The Confounding Problem in Existing Comparisons}
\label{sec:confounding}

A critical methodological limitation pervades the existing literature
on district heating modeling: published comparisons between different
topology representations invariably change the physics model and the
spatial aggregation level simultaneously, making it impossible to
attribute observed differences to either factor independently.

This confounding manifests through three co-varying dimensions:

\begin{enumerate}
    \item \textbf{Loss formulation}: Copperplate models typically use
    no losses or constant loss factors; pipe-level models employ
    temperature-dependent steady-state or dynamic loss equations.

    \item \textbf{COP treatment}: Single-node models often use
    constant or seasonally-averaged COP values; detailed models
    couple COP to network supply temperatures.

    \item \textbf{Temporal resolution}: Screening tools frequently
    operate at lower temporal resolution than detailed simulation
    models.
\end{enumerate}

To illustrate the problem concretely: Vesterlund et al.~\cite{Vesterlund2017}
use a full nonlinear pipe model with temperature-dependent COP, while
Gabrielli et al.~\cite{Gabrielli2018} apply a copperplate assumption
with constant thermal efficiencies. Any comparison of their respective
cost results would conflate topology effects with (1)~NLP vs.\
constant-loss formulations, (2)~temperature-dependent vs.\ constant
COP, and (3)~potentially differing temporal resolutions. The same
pattern appears when contrasting the zonal models of Wirtz et
al.~\cite{Wirtz2020} against full-graph representations in Schweiger
et al.~\cite{Schweiger2017}.

The consequence is that the practitioner question that matters
most---\emph{how much does topology abstraction itself change
operational cost and emission estimates, for a system with fixed
physics?}---remains unanswered.

The present work resolves this confounding through controlled
experimental design. Three topology abstraction levels (L1, L2, L3)
are embedded within a \emph{unified} MILP framework where the
physics model (loss formulation, COP parameterization, component
efficiencies) remains \emph{identical} at L2 and L3, while L1 omits
network physics as a deliberate baseline. Topology is thereby
isolated as the sole experimental variable.

\subsection{Research Questions}
\label{sec:rq}

\begin{description}[style=nextline]
    \item[RQ1 --- Topology Cost Impact:]
    How much does topology abstraction change annual operational cost
    and \ce{CO2} emissions when physics are held constant across levels?

    \item[RQ2 --- Dispatch Sensitivity:]
    Which dispatch decisions (heat pump utilization, storage charging
    patterns, boiler activation frequency) are most sensitive to the
    topology level, and what is the physical mechanism?

    \item[RQ3 --- Generalizability:]
    Which network properties (pipe length, demand heterogeneity,
    storage capacity) determine whether detailed topology
    representation is warranted?
\end{description}

\subsection{Contributions}
\label{sec:contributions}

\begin{enumerate}
    \item \textbf{L1/L2/L3 taxonomy}: A reusable, transferable
    classification framework for district heating MILP modeling
    fidelity, enabling consistent comparison across future studies
    (\Cref{sec:topology_def}).

    \item \textbf{First controlled comparison}: To the best of the
    authors' knowledge, no prior study has isolated topology as the
    sole experimental variable while holding physics constant across
    L2 and L3. This paper fills that gap through a comparison
    validated against measured industrial operational data
    (\Cref{sec:results}).

    \item \textbf{Storage dispatch as primary diagnostic}: Mechanistic
    explanation of how topology affects cost via loss-dependent
    temporal holding costs, with storage charge timing identified as
    a directly observable indicator of topology adequacy
    (\Cref{sec:results_storage}).

    \item \textbf{Actionable decision criteria}: Parametric analysis
    across 36 synthetic network configurations yielding quantitative
    topology-selection rules based on observable network properties
    (\Cref{sec:generalizability}).
\end{enumerate}

\subsection{Paper Structure}

\Cref{sec:related_work} reviews related work and positions this
contribution within the district heating optimization literature.
\Cref{sec:model} presents the unified MILP formulation.
\Cref{sec:case_studies} describes the primary industrial case study
and the synthetic parametric network.
\Cref{sec:results} presents primary results.
\Cref{sec:generalizability} extends findings through parametric
analysis. \Cref{sec:sensitivity} provides sensitivity analysis.
\Cref{sec:discussion} synthesizes findings.
\Cref{sec:conclusion} concludes with recommendations.

%% ============================================================
%% 2. RELATED WORK
%% ============================================================
\section{Related Work}
\label{sec:related_work}

\subsection{MILP Optimization of District Heating Systems}
\label{sec:rw_milp}

Mixed-integer linear programming has become the dominant methodology
for operational optimization and capacity planning of district heating
systems~\cite{Fang2016,Guelpa2019review}. The technology combination
of heat pumps, thermal storage, CHP units, and electrode boilers
that characterizes modern electrified networks introduces both binary
commitment decisions and nonlinear physical relationships that must
be handled through model approximation~\cite{Lund2014_4GDH}.

Several open-source frameworks support multi-carrier energy system
optimization. The Open Energy Modelling Framework
(oemof)~\cite{Hilpert2018_oemof} provides modular component libraries
including the oemof-thermal extension for thermal
networks~\cite{Krien2020_oemofthermal}. PyPSA~\cite{Brown2018_pypsa}
focuses on power systems with heat sector coupling via linearized
models. The FINE framework~\cite{Welder2018_FINE} supports
multi-regional energy system optimization with explicit inter-regional
transmission. Calliope~\cite{Pfenninger2018_calliope} offers flexible
spatial and temporal resolution. The urbs framework~\cite{Dorfner2016_urbs}
and Balmorel~\cite{Wiese2018_balmorel} further extend the landscape.

Despite their diversity, these frameworks share common limitations
for detailed district heating analysis. Thermal losses are typically
represented through aggregate, time-invariant loss factors rather than
physics-based temperature-dependent pipe models. Heat pump COP is
often constant or seasonal, breaking the physical coupling to network
supply temperatures. Network topology is either fully aggregated
(copperplate) or fixed at a single detail level without capability
for systematic level comparison. The present work addresses these
limitations within a specialized framework retaining MILP tractability.

\subsection{Topology Representation in District Heating Models}
\label{sec:rw_topology}

The district heating modeling literature exhibits a continuous
spectrum of spatial aggregation levels, which we organize into
three categories corresponding directly to our L1/L2/L3 taxonomy.

\textbf{Single-node aggregation (L1-equivalent):} All generation
and demand are aggregated to a single bus. This representation
is standard in national-scale energy system models where heating
is one of many sectors~\cite{Connolly2014_EnergyPLAN,Brown2018_pypsa,
Lund2017_SmartEnergy}. Spatial structure and thermal losses are absent.

\textbf{Zone-based models (L2-equivalent):} Demand is aggregated
into geographic zones (typically 3--15 nodes) connected by simplified
inter-zone pipe segments with aggregate pipe parameters~\cite{Wirtz2020,
Mertz2016,Buffa2019}. Reduced-order models by Bünning et
al.~\cite{Buenning2020} demonstrate that zone-level aggregation
captures 85--95\% of dynamic temperature variation in typical networks.

\textbf{Full-graph models (L3-equivalent):} Node-by-node
representation matching the physical network topology, with
pipe-by-pipe parameters sourced from GIS data~\cite{Vesterlund2017,
Schweiger2017,vanderHeijde2017,Blommaert2020}. Giraud et
al.~\cite{Giraud2015} demonstrate that full-graph MILP formulations
for branching networks remain tractable up to approximately 100~nodes
for hourly, annual optimization horizons.

\Cref{tab:prior_work} positions representative studies from each
category. The critical observation is that in all prior comparative
work, the physics model changes simultaneously with the topology
level.

\begin{table*}[htbp]
\centering
\caption{Representative district heating optimization studies by
topology level, physics model, and whether topology effects are
isolated from physics effects.}
\label{tab:prior_work}
\small
\begin{tabular}{@{}p{3.5cm}p{2.0cm}p{3.0cm}p{2.2cm}p{1.5cm}p{1.8cm}@{}}
\toprule
Reference & Topology & Physics Model & COP Treatment & Scale & Controlled? \\
\midrule
Connolly et al.~\cite{Connolly2014_EnergyPLAN}
  & Copperplate & Fixed efficiency & Constant seasonal & National & No \\
Gabrielli et al.~\cite{Gabrielli2018}
  & Copperplate & Fixed efficiency & Constant & 10~MW & No \\
Mertz et al.~\cite{Mertz2016}
  & Zone (8 nodes) & Linear losses & Constant & 20~MW & No \\
Wirtz et al.~\cite{Wirtz2020}
  & Zone (5 nodes) & Linear losses & Fixed seasonal & 25~MW & No \\
Bünning et al.~\cite{Buenning2020}
  & Reduced-order & Simplified & Constant & 8~MW & No \\
Giraud et al.~\cite{Giraud2015}
  & Full graph & Steady-state & Temperature-dep. & 30~MW & No \\
Vesterlund et al.~\cite{Vesterlund2017}
  & Full graph & NLP thermal & Temperature-dep. & 50~MW & No \\
Schweiger et al.~\cite{Schweiger2017}
  & Full graph & Modelica co-sim & Dynamic & 15~MW & No \\
van der Heijde et al.~\cite{vanderHeijde2017}
  & Full graph & Pipe model comp. & -- & Review & Partial$^{a}$ \\
Blommaert et al.~\cite{Blommaert2020}
  & Full graph & Adjoint NLP & Temperature-dep. & 30~MW & No \\
\textbf{This work}
  & \textbf{L1/L2/L3} & \textbf{Identical MILP} & \textbf{Pre-computed}
  & \textbf{$\approx$12~MW} & \textbf{Yes} \\
\bottomrule
\multicolumn{6}{@{}l}{\footnotesize $^{a}$~Compared pipe loss
formulations but not spatial aggregation levels.}
\end{tabular}
\end{table*}

\subsection{Thermal Pipe Loss Modeling}
\label{sec:rw_losses}

Pipe thermal losses follow from steady-state or dynamic heat transfer
between the fluid and surrounding soil. The steady-state linear model,
in which loss is proportional to the temperature difference between
fluid and ground~\cite{Paalsson1999,Frederiksen2013}, forms the basis
for MILP-compatible loss calculations used in this study.
Dynamic pipe models capture transport delays and thermal inertia but
introduce nonlinearity or require finite-element
discretization~\cite{vanderHeijde2017,Duquette2016}.

The significance of network losses as an optimization variable depends
on their magnitude relative to demand. Werner~\cite{Werner2017} reports
European DH network heat loss ratios of 5--25\%, with modern
pre-insulated networks typically below 10\%. For systems in this
range, ignoring losses (L1) introduces a proportional underestimation
of generation requirements and cost. Talebi et al.~\cite{Talebi2016}
confirm that loss fraction increases with supply temperature and pipe
length, motivating the systematic parameter variation in the synthetic
analysis.

\subsection{Thermal Storage Dispatch in District Heating}
\label{sec:rw_storage}

Thermal energy storage (TES) provides flexibility in district heating
through peak shaving, load shifting, and renewable
integration~\cite{Vandermeulen2018,Guelpa2019review}. Christidis et
al.~\cite{Christidis2017} demonstrate that optimal TES dispatch in
CHP-based networks shifts heat production by 4--8~hours to exploit
overnight electricity pricing---a pattern directly dependent on
temporal cost signals.

The economic value of storage through temporal arbitrage depends on
the accuracy of the cost model used to compute holding costs. In
copperplate models, the only time-varying cost is the electricity
purchase price. In network-aware models, an additional implicit cost
appears: heat discharged during high-demand winter periods must
compensate for losses that are also highest in those periods (large
$\Delta T$ between fluid and cold ground). This asymmetric temporal
cost structure modifies optimal charging decisions in ways invisible
to L1 models. The present paper quantifies this mechanism for the
first time in a controlled setting.

\subsection{Research Gap}
\label{sec:rw_gap}

Synthesizing the above review, three gaps remain unaddressed:

\begin{enumerate}
    \item No study isolates topology abstraction as the sole
    experimental variable while holding physics constant. All
    existing cross-level comparisons confound topology with physics.

    \item No quantitative guidance exists for topology selection
    as a function of observable network properties (pipe length,
    demand heterogeneity, storage capacity).

    \item The mechanism by which topology level affects storage
    dispatch---through loss-dependent temporal holding
    costs---has not been analyzed or quantified.
\end{enumerate}

The present work addresses all three gaps within a single controlled
study.

%% ============================================================
%% 3. MODEL FORMULATION
%% ============================================================
\section{Model Formulation}
\label{sec:model}

\subsection{Experimental Design: Topology as Sole Variable}
\label{sec:exp_design}

The central methodological principle is controlled isolation of
topology effects. The physics model is deliberately simplified
and---critically---kept \emph{identical} across L2 and L3:

\begin{itemize}
    \item \textbf{Steady-state thermal losses}: Physics-based per-pipe
    losses computed separately for supply and return pipes
    (\Cref{sec:model_losses}).

    \item \textbf{Fixed supply temperature}: Determined by the
    heating curve $T_{\mathrm{sup}}^{\mathrm{nom}}(t)$, not optimized
    as a decision variable. This eliminates temperature--flow coupling
    while preserving seasonally-varying loss magnitudes.

    \item \textbf{Pre-computed COP}: Time-varying parameters from
    bilinear interpolation of manufacturer lookup tables
    (\Cref{sec:model_hp}). COP variation is captured but enters the
    model as a parameter, preserving MILP linearity.

    \item \textbf{No transport delays}: Heat propagation is assumed
    instantaneous within each hourly timestep. For the pipe lengths
    and flow velocities of the case study network, transport times
    are estimated at 0.2--1.8~h. Introducing delays would confound
    topology with temporal dynamics.

    \item \textbf{No pressure optimization}: Hydraulic pressure drop
    is excluded from the optimization. Post-processing via the
    Darcy--Weisbach equation at peak flow conditions confirms that
    pumping power is $<$2\% of total network thermal losses.

    \item \textbf{Dispatch-only horizon}: All component capacities
    fixed; no investment decisions.
\end{itemize}

These are \emph{design choices for experimental control}, not
framework limitations. The underlying optimization framework supports
extended physics including SOS2-linearized transport delays,
piecewise-linear pressure drop coupling, and unit commitment
constraints. These capabilities are documented in a companion
paper~\cite{Paper2_CALION} and are deactivated for this controlled
study.

\subsection{L1/L2/L3 Topology Definitions}
\label{sec:topology_def}

Three topology levels span the practical range of district heating
representations.

\textbf{Level~1 (Copperplate --- L1):} Single-node representation
with no spatial structure. All generators, storage units, and demands
connect to a single aggregated bus. No pipe losses: $Q_{\mathrm{loss}}=0$
for all $t$. The heat balance is a single scalar equation.

\textbf{Level~2 (Simplified multi-node --- L2):} Multiple aggregated
demand zones (5 nodes in the Stadtbach case) connected by simplified
pipe routing. Physics-based steady-state thermal losses computed per
inter-zone pipe segment using aggregated parameters. The total
$U \cdot L$ product is preserved from L3 (see \Cref{sec:case_stadtbach}),
ensuring that aggregate annual loss energy is matched between L2 and
L3 under identical temperature conditions. Only the spatial distribution
of losses differs.

\textbf{Level~3 (Detailed multi-node --- L3):} Full spatial resolution
matching the physical network geometry from GIS data (30 nodes in the
Stadtbach case). Individual pipe parameters. Same steady-state thermal
loss formulation as L2.

\textbf{Critical design constraint:} L2 and L3 use identical thermal
loss \emph{formulations}. Total pipe length per type and aggregate
$U$-values are matched between levels. Thermal storage capacity and
component installed capacities are identical. Only the number of nodes
and the pipe-level spatial resolution differ.

\Cref{tab:feature_matrix} specifies which model features are active
at each level.

\begin{table}[htbp]
\centering
\caption{Feature activation matrix across topology levels. All
features active at L2/L3 use identical formulations.}
\label{tab:feature_matrix}
\begin{tabular}{@{}lccc@{}}
\toprule
Model Feature & L1 & L2 & L3 \\
\midrule
Aggregated heat balance, Eq.~\eqref{eq:balance_L1}
  & \checkmark & -- & -- \\
Nodal heat balances, Eq.~\eqref{eq:balance_prod}--\eqref{eq:balance_consumer}
  & -- & \checkmark & \checkmark \\
Pipe thermal losses, Eq.~\eqref{eq:loss_supply}--\eqref{eq:loss_return}
  & -- & \checkmark & \checkmark \\
Global electricity balance, Eq.~\eqref{eq:balance_el}
  & \checkmark & \checkmark & \checkmark \\
Heat pump (dual-source), Eq.~\eqref{eq:hp_split}--\eqref{eq:hp_cap}
  & \checkmark & \checkmark & \checkmark \\
Thermal storage (SOC dynamics), Eq.~\eqref{eq:storage_soc}
  & \checkmark & \checkmark & \checkmark \\
Gas boiler \& electrode boiler, Eq.~\eqref{eq:gen_thermal},\,\eqref{eq:p2h}
  & \checkmark & \checkmark & \checkmark \\
Grid buy/sell logic (Big-M), Eq.~\eqref{eq:grid_bigm}
  & \checkmark & \checkmark & \checkmark \\
Peak import demand charge, Eq.~\eqref{eq:peak}
  & \checkmark & \checkmark & \checkmark \\
Hourly \ce{CO2} accounting, Eq.~\eqref{eq:em_grid}--\eqref{eq:em_fuel}
  & \checkmark & \checkmark & \checkmark \\
\midrule
\textit{Disabled (experimental control):} & & & \\
Transport delays (flow-dependent) & -- & -- & -- \\
Pressure drop optimization & -- & -- & -- \\
Unit commitment (min up/down time) & -- & -- & -- \\
Investment sizing decisions & -- & -- & -- \\
\bottomrule
\end{tabular}
\end{table}

\subsection{Objective Function}
\label{sec:model_objective}

The model minimizes total annual operational cost:
\begin{equation}
\min \; Z = C^{\mathrm{en}} + C^{\mathrm{fuel}}
+ C^{\mathrm{dump}} + C^{\ce{CO2}} + C^{\mathrm{dem}}
\label{eq:objective}
\end{equation}

\textbf{Electricity procurement and sales:}
\begin{equation}
C^{\mathrm{en}} = \Delta t \sum_{t \in \mathcal{T}} \left(
P_t^{\mathrm{buy}} \lambda_t^{\mathrm{buy}}
- P_t^{\mathrm{sell}} \lambda_t^{\mathrm{sell}} \right)
\label{eq:cost_energy}
\end{equation}

The effective buying price includes market price, grid fees, and levies:
\begin{equation}
\lambda_t^{\mathrm{buy}} = \lambda_t^{\mathrm{spot}}
+ \lambda^{\mathrm{fee}} + \lambda^{\mathrm{grid}}
\label{eq:buy_price}
\end{equation}

The selling price accounts for bid-ask spread, floor, haircut,
export fee, and balancing premium (full derivation in
\Cref{app:selling_price}):
\begin{equation}
\lambda_t^{\mathrm{sell}} = \max\!\Big(\!
\max\!\big(\lambda_t^{\mathrm{spot}} - \lambda^{\mathrm{spread}},\,
\lambda^{\mathrm{floor}}\big) (1{-}h^{\mathrm{sell}})
- \lambda^{\mathrm{sell,fee}} + \lambda^{\mathrm{sell,prem}},\; 0\Big)
\label{eq:sell_price}
\end{equation}

\textbf{Fuel cost:}
\begin{equation}
C^{\mathrm{fuel}} = \Delta t \sum_{g \in \mathcal{G}^{\mathrm{fuel}}}
\sum_{t \in \mathcal{T}} F_{g,t} \lambda_g^{\mathrm{fuel}}
\label{eq:cost_fuel}
\end{equation}

\textbf{Heat curtailment penalty:}
\begin{equation}
C^{\mathrm{dump}} = \Delta t \sum_{n \in \mathcal{N}}
\sum_{t \in \mathcal{T}} Q_{n,t}^{\mathrm{dump}} \lambda^{\mathrm{dump}}
\label{eq:cost_dump}
\end{equation}

\textbf{\ce{CO2} cost:}
\begin{equation}
C^{\ce{CO2}} = \frac{\lambda^{\ce{CO2}}}{1000} \left(
\sum_{t} E_t^{\ce{CO2},\mathrm{grid}}
+ \sum_{g,t} E_{g,t}^{\ce{CO2},\mathrm{fuel}} \right)
\label{eq:cost_co2}
\end{equation}

\textbf{Demand charge} (annual peak import fee, prorated by year
fraction):
\begin{equation}
C^{\mathrm{dem}} = \lambda^{\mathrm{dem}} \cdot f^{\mathrm{yr}}
\cdot P^{\mathrm{peak}}
\label{eq:cost_demand}
\end{equation}

\subsection{Energy Balance Constraints}
\label{sec:model_balance}

\subsubsection{Thermal Energy Balances}

\textbf{L1 --- Aggregated copperplate balance:}
\begin{equation}
\sum_{g \in \mathcal{G}^{\mathrm{th}}} Q_{g,t}^{\mathrm{th}}
+ \sum_{s \in \mathcal{S}} Q_{s,t}^{\mathrm{dis}}
= Q_t^{\mathrm{dem}} + Q_t^{\mathrm{dump}}
+ \sum_{s \in \mathcal{S}} Q_{s,t}^{\mathrm{ch}}
\qquad \forall t
\label{eq:balance_L1}
\end{equation}

L1 carries no loss term. The total demand $Q_t^{\mathrm{dem}}$
aggregates all nodal demands. This lossless balance is the defining
characteristic of L1.

\textbf{L2/L3 --- Nodal balance, producer nodes:}
\begin{equation}
\sum_{g \in \mathcal{G}_n^{\mathrm{th}}} Q_{g,t}^{\mathrm{th}}
+ \sum_{s \in \mathcal{S}_n} Q_{s,t}^{\mathrm{dis}}
= Q_{n,t}^{\mathrm{dem}} + Q_{n,t}^{\mathrm{dump}}
+ \sum_{s \in \mathcal{S}_n} Q_{s,t}^{\mathrm{ch}}
+ Q_{n,t}^{\mathrm{loss}}
\qquad \forall n \in \mathcal{N}^{\mathrm{prod}},\, \forall t
\label{eq:balance_prod}
\end{equation}

\textbf{L2/L3 --- Nodal balance, consumer nodes:}
\begin{equation}
Q_{n,t}^{\mathrm{pipe}}
+ \sum_{g \in \mathcal{G}_n^{\mathrm{th}}} Q_{g,t}^{\mathrm{th}}
+ \sum_{s \in \mathcal{S}_n} Q_{s,t}^{\mathrm{dis}}
= Q_{n,t}^{\mathrm{dem}} + Q_{n,t}^{\mathrm{dump}}
+ \sum_{s \in \mathcal{S}_n} Q_{s,t}^{\mathrm{ch}}
\qquad \forall n \in \mathcal{N}^{\mathrm{con}},\, \forall t
\label{eq:balance_consumer}
\end{equation}

Here $Q_{n,t}^{\mathrm{pipe}}$ is heat delivered to consumer node $n$
via the connecting pipe (after loss deduction), and $Q_{n,t}^{\mathrm{loss}}$
is the total thermal loss allocated to producer node $n$.

Total network losses:
\begin{equation}
Q_t^{\mathrm{loss,total}} = \sum_{p \in \mathcal{P}} \left(
Q_{p,t}^{\mathrm{loss,sup}} + Q_{p,t}^{\mathrm{loss,ret}} \right)
\qquad \forall t
\label{eq:loss_total}
\end{equation}

\subsubsection{Global Electricity Balance}

\begin{equation}
P_t^{\mathrm{buy}} + \sum_{g \in \mathcal{G}^{\mathrm{CHP}}}
P_{g,t}^{\mathrm{el}}
= \sum_{h \in \mathcal{H}} P_{h,t}^{\mathrm{el}}
+ \sum_{r \in \mathcal{R}} P_{r,t}^{\mathrm{el}}
+ P_t^{\mathrm{sell}}
\qquad \forall t
\label{eq:balance_el}
\end{equation}

Mutual exclusivity of grid import and export:
\begin{equation}
P_t^{\mathrm{buy}} \leq M^{\mathrm{grid}} z_t, \qquad
P_t^{\mathrm{sell}} \leq M^{\mathrm{grid}} (1 - z_t)
\qquad \forall t
\label{eq:grid_bigm}
\end{equation}

Grid connection capacity limits:
\begin{equation}
P_t^{\mathrm{buy}} \leq \bar{P}^{\mathrm{imp}}, \qquad
P_t^{\mathrm{sell}} \leq \bar{P}^{\mathrm{exp}}
\qquad \forall t
\label{eq:grid_cap}
\end{equation}

Peak import tracking for demand charge computation:
\begin{equation}
P^{\mathrm{peak}} \geq P_t^{\mathrm{buy}} \qquad \forall t
\label{eq:peak}
\end{equation}

\subsection{Thermal Loss Model (L2/L3 Only)}
\label{sec:model_losses}

Thermal losses follow from steady-state conduction, computed
\emph{separately} for supply and return pipes:

\textbf{Supply pipe losses:}
\begin{equation}
Q_{p,t}^{\mathrm{loss,sup}} = \frac{U_p \cdot L_p \cdot
\big(T_{\mathrm{sup}}^{\mathrm{nom}}(t) - T_{\mathrm{ground}}(t)\big)}{10^6}
\quad [\mathrm{MW}]
\label{eq:loss_supply}
\end{equation}

\textbf{Return pipe losses:}
\begin{equation}
Q_{p,t}^{\mathrm{loss,ret}} = \frac{U_p \cdot L_p \cdot
\big(T_{\mathrm{ret}}^{\mathrm{nom}} - T_{\mathrm{ground}}(t)\big)}{10^6}
\quad [\mathrm{MW}]
\label{eq:loss_return}
\end{equation}

Here $U_p$~[\si{W/(m.K)}] is the linear heat transfer coefficient
(accounting for insulation, soil, and pipe wall), $L_p$~[m] the pipe
length, and $T_{\mathrm{ground}}(t)$ the measured ground temperature at
installation depth. The factor $10^{-6}$ converts from \si{W} to \si{MW}.

Losses are time-varying through $T_{\mathrm{ground}}(t)$ and
$T_{\mathrm{sup}}^{\mathrm{nom}}(t)$, with seasonal amplitude of
approximately \placeholder{X}~MW\textsubscript{th} in the case study
network.

\textbf{L1 treatment:} $Q_{p,t}^{\mathrm{loss,sup}} =
Q_{p,t}^{\mathrm{loss,ret}} = 0$ for all $p, t$. The complete
omission of losses constitutes the primary physical difference
between L1 and L2/L3.

\textbf{Validity of steady-state assumption:} The steady-state
approximation is accurate when the thermal Peclet number
$Pe = \dot{m} c_p L / (U A)$ is large~\cite{vanderHeijde2017}.
For the case study network at average flow velocities (0.5--1.5~m/s),
$Pe$ ranges from \placeholder{XX--XX}, confirming that advective heat
transport dominates within the one-hour timestep. Dynamic effects
manifest primarily in temperature transients during startup, which are
outside the scope of the hourly dispatch model.

\subsection{Heat Pump Model}
\label{sec:model_hp}

Each heat pump $h$ supports dual heat source operation: a
waste-heat-recovery (WRG) source with time-varying COP, and a default
source (ambient air or ground) with constant COP.

Total heat output:
\begin{equation}
Q_{h,t} = Q_{h,t}^{\mathrm{wrg}} + Q_{h,t}^{\mathrm{def}}
\qquad \forall h,\, \forall t
\label{eq:hp_split}
\end{equation}

Electrical power consumption:
\begin{equation}
P_{h,t}^{\mathrm{el}} = \frac{Q_{h,t}^{\mathrm{wrg}}}{
\mathrm{COP}_{h,t}^{\mathrm{wrg}}}
+ \frac{Q_{h,t}^{\mathrm{def}}}{\mathrm{COP}_h^{\mathrm{def}}}
\qquad \forall h,\, \forall t
\label{eq:hp_pel}
\end{equation}

Both COP values enter as parameters, preserving MILP linearity.
$\mathrm{COP}_{h,t}^{\mathrm{wrg}}$ is pre-computed via bilinear
interpolation on the manufacturer performance map:
\begin{equation}
\mathrm{COP}_{h,t}^{\mathrm{wrg}} = f_{\mathrm{bilinear}}\big(
T_{\mathrm{source}}(t),\, T_{\mathrm{sink}}(t)\big),
\qquad \mathrm{COP} \in [0.5,\, 15.0]
\label{eq:cop_lookup}
\end{equation}

The clamping bounds $[0.5, 15.0]$ prevent numerical artifacts at
extreme temperature differentials; they activate in $<$0.3\% of
timesteps. Interpolation error is bounded at $\pm$2--3\% relative
COP (derivation in \Cref{app:cop}), translating to $\pm$1--2\%
system cost variation.

Capacity and minimum part-load constraints:
\begin{equation}
\phi_h^{\mathrm{min}} \bar{Q}_h u_{h,t} \leq Q_{h,t}
\leq \bar{Q}_h u_{h,t}
\qquad \forall h,\, \forall t
\label{eq:hp_cap}
\end{equation}

\subsection{Thermal Storage Model}
\label{sec:model_storage}

State-of-charge dynamics with multiplicative self-discharge
(thermodynamically consistent for variable $\Delta t$):
\begin{equation}
E_{s,t} = (1-\alpha_s)^{\Delta t} \cdot E_{s,t-1}
+ \eta_{s}^{\mathrm{ch}} Q_{s,t}^{\mathrm{ch}} \Delta t
- \frac{Q_{s,t}^{\mathrm{dis}} \Delta t}{\eta_{s}^{\mathrm{dis}}}
\qquad \forall s,\, \forall t
\label{eq:storage_soc}
\end{equation}

Energy and power capacity constraints:
\begin{equation}
E_{s,t} \leq \bar{E}_s \cdot a_{s,t}
\qquad \forall s,\, \forall t
\label{eq:storage_cap_E}
\end{equation}
\begin{equation}
Q_{s,t}^{\mathrm{ch}} \leq \bar{P}_s \delta_{s,t}^{\mathrm{ch}},
\qquad
Q_{s,t}^{\mathrm{dis}} \leq \bar{P}_s \delta_{s,t}^{\mathrm{dis}}
\qquad \forall s,\, \forall t
\label{eq:storage_cap_P}
\end{equation}

Simultaneous charging and discharging prevention:
\begin{equation}
\delta_{s,t}^{\mathrm{ch}} + \delta_{s,t}^{\mathrm{dis}}
\leq a_{s,t}
\qquad \forall s,\, \forall t
\label{eq:storage_excl}
\end{equation}

\textbf{Why storage is the primary diagnostic variable:} The temporal
arbitrage value of storage depends on holding costs. In L1, no
physical holding cost exists beyond self-discharge---charging responds
purely to electricity price differentials. In L2/L3, network losses
introduce an implicit time-varying holding cost: heat stored at time
$t_1$ for discharge at $t_2$ must additionally compensate for the loss
overhead at $t_2$. The effective arbitrage spread becomes
$\lambda_{t_1}^{\mathrm{buy}} - (\lambda_{t_2}^{\mathrm{buy}} +
\lambda_{t_2}^{\mathrm{loss,marginal}})$, where
$\lambda_{t_2}^{\mathrm{loss,marginal}}$ is the marginal generation
cost of covering the loss increment at $t_2$. During cold winter
afternoons this penalty is highest, discouraging the deep overnight
charging that L1 would otherwise schedule.
\Cref{sec:results_storage} quantifies the resulting shift in charging
hour distributions.

\subsection{Thermal Generator Models}
\label{sec:model_gen}

\textbf{Gas boiler:}
\begin{equation}
Q_{g,t}^{\mathrm{th}} = \eta_g^{\mathrm{th}} F_{g,t},
\qquad
Q_{g,t}^{\mathrm{th}} \leq \bar{Q}_g^{\mathrm{th}}
\qquad \forall g \in \mathcal{G}^{\mathrm{th}},\, \forall t
\label{eq:gen_thermal}
\end{equation}

\textbf{CHP unit:}
\begin{equation}
Q_{g,t}^{\mathrm{th}} = \eta_g^{\mathrm{th}} F_{g,t}, \qquad
P_{g,t}^{\mathrm{el}} = \eta_g^{\mathrm{el}} F_{g,t}
\qquad \forall g \in \mathcal{G}^{\mathrm{CHP}},\, \forall t
\label{eq:gen_chp}
\end{equation}

\textbf{Electrode boiler / Power-to-Heat:}
\begin{equation}
Q_{r,t} = \eta_{r} P_{r,t}^{\mathrm{el}}, \qquad
Q_{r,t} \leq \bar{Q}_r
\qquad \forall r \in \mathcal{R},\, \forall t
\label{eq:p2h}
\end{equation}

\subsection{Emissions Accounting}
\label{sec:model_emissions}

Hourly grid electricity emissions:
\begin{equation}
E_t^{\ce{CO2},\mathrm{grid}} = P_t^{\mathrm{buy}}
\cdot e_t^{\mathrm{grid}} \cdot \Delta t
\qquad \forall t
\label{eq:em_grid}
\end{equation}

Fuel emissions:
\begin{equation}
E_{g,t}^{\ce{CO2},\mathrm{fuel}} = F_{g,t} \cdot e_g^{\mathrm{fuel}}
\cdot \Delta t
\qquad \forall g,\, \forall t
\label{eq:em_fuel}
\end{equation}

For CHP units, emissions are allocated between heat and electricity
outputs using the efficiency-weighted method~\cite{Lund2014_4GDH},
chosen for consistency with the Austrian national reporting framework:
\begin{equation}
\alpha_g^{\mathrm{heat}} = \frac{\eta_g^{\mathrm{th}}}{
\eta_g^{\mathrm{th}} + \eta_g^{\mathrm{el}}}, \qquad
\alpha_g^{\mathrm{el}} = \frac{\eta_g^{\mathrm{el}}}{
\eta_g^{\mathrm{th}} + \eta_g^{\mathrm{el}}}
\label{eq:em_chp_alloc}
\end{equation}

\subsection{Demand Heterogeneity Index}
\label{sec:model_hi}

The demand heterogeneity index HI quantifies the spatial
non-uniformity of annual nodal demand shares:
\begin{equation}
d_n = \frac{\sum_t Q_{n,t}^{\mathrm{dem}}}{\sum_n \sum_t
Q_{n,t}^{\mathrm{dem}}}, \qquad
\mathrm{HI} = \frac{N}{N-1} \sum_{n \in \mathcal{N}}
\left(d_n - \frac{1}{N}\right)^2
\label{eq:hi}
\end{equation}

where $d_n$ is node $n$'s fraction of total annual demand and
$N$ the node count. HI~$= 0$ corresponds to perfectly uniform demand;
HI~$= 1$ to complete concentration at a single node. The
normalization factor $N/(N-1)$ corrects for the bias of the sample
variance estimator and ensures comparability across networks with
different node counts.

\subsection{Implementation}
\label{sec:model_impl}

The model is implemented in Python~3.10 using the Pyomo algebraic
modeling language~\cite{pyomo}. Network configuration (topology,
component parameters, boundary conditions) is specified via YAML
files enabling L1/L2/L3 instantiation from a single parameterized
codebase. Solver: Gurobi~10.0 with MIP gap tolerance 0.01\%, ensuring
that reported cost differences $\geq$0.1\% are attributable to model
structure rather than solver tolerance. All MILP problems use identical
solver settings and random seeds.

%% ============================================================
%% 4. CASE STUDIES
%% ============================================================
\section{Case Studies}
\label{sec:case_studies}

\subsection{Primary Case Study: Industrial Heat Network (Stadtbach)}
\label{sec:case_stadtbach}

The primary validation case is an Austrian industrial heat network
providing process heat to a multi-building site with connection to
the public electricity grid. The network integrates multiple
electrified heat supply technologies and thermal storage, representing
the increasingly common configuration of industrial 4GDH
installations~\cite{Lund2014_4GDH}.

\begin{table}[htbp]
\centering
\caption{Stadtbach network: installed components and system parameters.
Specific values anonymized per agreement with the industrial partner.}
\label{tab:stadtbach}
\begin{tabular}{@{}lll@{}}
\toprule
Component / Parameter & Value & Notes \\
\midrule
\multicolumn{3}{@{}l}{\textit{Demand}} \\
Peak heat demand           & \placeholder{8--15}~MW        & Industrial process heat \\
Annual heat demand         & \placeholder{25--50}~GWh      & \\
\midrule
\multicolumn{3}{@{}l}{\textit{Network}} \\
Total pipe length (supply+return) & \placeholder{3--8}~km  & \\
Demand nodes (L3)          & 30                            & From GIS data \\
Demand zones (L2)          & 5                             & N/S/E/W/Central \\
Demand heterogeneity (HI)  & \placeholder{0.XX}            & Eq.~\eqref{eq:hi} \\
Design supply temperature  & \placeholder{XX}~\si{\celsius} & \\
Design return temperature  & \placeholder{XX}~\si{\celsius} & \\
\midrule
\multicolumn{3}{@{}l}{\textit{Generation components}} \\
Heat pump (air-source, dual-source) & \placeholder{3--6}~MW\textsubscript{th} & COP 2.5--4.0 \\
Gas boiler                 & \placeholder{8--12}~MW\textsubscript{th} & Backup and peak \\
Electrode boiler (P2H)     & \placeholder{1--3}~MW\textsubscript{th}  & $\eta = 0.99$ \\
\midrule
\multicolumn{3}{@{}l}{\textit{Flexibility components}} \\
Thermal storage            & \placeholder{50--150}~MWh     & Stratified hot water \\
Storage charging ratio     & \placeholder{X.X}~h of peak   & \\
\midrule
\multicolumn{3}{@{}l}{\textit{Grid coupling}} \\
Import capacity            & \placeholder{XX}~MW           & Transformer limit \\
Export capacity            & \placeholder{XX}~MW           & Bilateral contract \\
\bottomrule
\end{tabular}
\end{table}

\subsection{L1/L2/L3 Implementations of the Stadtbach Network}
\label{sec:case_levels}

\textbf{L3 (30 nodes):} Full physical topology from GIS data. Each
demand connection point is a distinct node. Individual pipe parameters
($U_p$, $L_p$) from network documentation. Total pipe length and
aggregate $U$-values serve as the reference matched by L2.

\textbf{L2 (5 nodes):} Geographic aggregation into five demand zones
(North, South, East, West, Central production hub). The total
$U \cdot L$ product per pipe type is preserved:
$\sum_{p \in \mathcal{P}^{\mathrm{L2}}} U_p L_p =
\sum_{p \in \mathcal{P}^{\mathrm{L3}}} U_p L_p$.
Demand at each L2 node equals the sum of constituent L3 nodal demands
at each timestep.

\textbf{L1 (1 node):} Single aggregated bus. Total demand equals the
sum of all 30 nodal demands. No losses.

\textbf{Controlled matching across levels:} Total installed capacity
of each component type, total pipe length, aggregate $U$-values, and
thermal storage parameters are identical in L2 and L3. All boundary
condition time series (electricity prices, temperatures, demand
profiles) are identical across all three levels.

\subsection{Input Data}
\label{sec:case_data}

\begin{itemize}
    \item \textbf{Electricity prices}: EPEX Austria day-ahead spot,
    hourly, 2022/2023 period, characterized by high price volatility
    and negative-price episodes.

    \item \textbf{Temperatures}: Site-measured ambient and ground
    temperature (at installation depth), hourly, used to compute
    $\mathrm{COP}_{h,t}$ and thermal losses simultaneously.

    \item \textbf{Heat demand}: Measured hourly at meter level.
    Disaggregated to L3 nodes, aggregated to L2 zones, and summed
    for L1. The same total demand applies to all levels.

    \item \textbf{Grid \ce{CO2} intensity}: ENTSO-E Transparency
    Platform, Austrian hourly generation-based emission factors.

    \item \textbf{Fuel prices}: Austrian wholesale natural gas
    (TTF-indexed), monthly average applied as constant within month.
\end{itemize}

\subsection{Validation Protocol}
\label{sec:case_validation}

L3 is validated against measured operational data \emph{prior to}
the controlled comparison using historical records including gas meter
readings, electricity import/export metering, heat pump running-hour
logs, and storage level measurements:

\begin{itemize}
    \item Annual gas consumption: MAPE target $<$5\%
    \item Annual electricity purchase: MAPE target $<$5\%
    \item Monthly heat pump operating hours: MAPE target $<$10\%
    \item Annual storage full-cycle equivalent: MAPE target $<$15\%
    \item Annual peak heat demand: MAPE target $<$5\%
\end{itemize}

Only L3 is validated against measured data. L1 and L2 are evaluated
\emph{relative to L3} as the validated reference. This is
methodologically consistent: the question is not whether L1 or L2
independently match reality, but how much they deviate from the
highest-fidelity model with identical physics.

\subsection{Synthetic Parametric Network}
\label{sec:case_synthetic}

A single real-world case study cannot establish generalizability. A
synthetic tree-topology network is parameterized by four independent
properties (\Cref{tab:synthetic}).

\begin{table}[htbp]
\centering
\caption{Synthetic network parameter space. Parameter ranges based on
the European district heating inventory~\cite{Werner2017,Frederiksen2013}.
All 81 combinations of three levels across four parameters are tested;
36 feasible configurations are retained after removing cases where pipe
length per node falls below the minimum physical separation distance
($<$50~m/node) or where storage ratio exceeds the maximum practical
charge/discharge cycle ($>$72~h).}
\label{tab:synthetic}
\begin{tabular}{@{}llllp{4.5cm}@{}}
\toprule
Parameter & Low & Medium & High & Physical rationale \\
\midrule
Total pipe length [km] & 1 & 5 & 15
  & Compact campus to extended urban district \\
Demand heterogeneity HI & 0.1 & 0.4 & 0.8
  & Near-uniform to strongly concentrated \\
Node count & 5 & 15 & 30
  & Sparse to dense spatial discretization \\
Storage ratio [h of peak] & 2 & 6 & 12
  & Small buffer to large shift capacity \\
\bottomrule
\end{tabular}
\end{table}

A tree topology is used because the Stadtbach network and the majority
of small- to medium industrial DH installations are radially structured.
Ring topologies introduce additional flow direction binary variables and
represent a distinct problem class; their treatment is deferred to future
work. The full experiment comprises 36 configurations $\times$ 3
topology levels = 108 MILP solves. Synthetic network geometry, demand
profiles, component parameters, and all configuration files are
published on Zenodo~\cite{Zenodo_data}.

%% ============================================================
%% 5. RESULTS
%% ============================================================
\section{Results: Stadtbach Case Study}
\label{sec:results}

\subsection{Model Validation (L3)}
\label{sec:results_validation}

\begin{table}[htbp]
\centering
\caption{L3 model validation against measured operational data.}
\label{tab:validation}
\begin{tabular}{@{}lllll@{}}
\toprule
Variable & Measured & L3 Model & MAPE & Target \\
\midrule
Annual gas consumption [MWh]      & \placeholder{XX,XXX} & \placeholder{XX,XXX} & \placeholder{X.X\%} & $<$5\% \\
Annual electricity purchase [MWh] & \placeholder{XX,XXX} & \placeholder{XX,XXX} & \placeholder{X.X\%} & $<$5\% \\
Annual HP operating hours [h]     & \placeholder{X,XXX}  & \placeholder{X,XXX}  & \placeholder{X.X\%} & $<$10\% \\
Storage full-cycle equivalent [--] & \placeholder{XXX}   & \placeholder{XXX}    & \placeholder{X.X\%} & $<$15\% \\
Peak heat demand [MW]             & \placeholder{XX}     & \placeholder{XX}     & \placeholder{X.X\%} & $<$5\% \\
\bottomrule
\end{tabular}
\end{table}

\subsection{Annual Operational Cost Comparison}
\label{sec:results_cost}

\begin{table*}[htbp]
\centering
\caption{Annual operational cost comparison across topology levels.
L3 is the validated reference. Negative $\Delta$ indicates
underestimation relative to L3.}
\label{tab:cost}
\begin{tabular}{@{}lllllll@{}}
\toprule
Level & Total [\euro/a] & $\Delta$ vs.\ L3 & Fuel [\euro/a]
  & Electricity [\euro/a] & \ce{CO2} [\euro/a] & Demand [\euro/a] \\
\midrule
L1 & \placeholder{X,XXX,XXX} & $-$\placeholder{X.X}\%
  & \placeholder{XXX,XXX} & \placeholder{XXX,XXX}
  & \placeholder{XX,XXX} & \placeholder{XX,XXX} \\
L2 & \placeholder{X,XXX,XXX} & $-$\placeholder{X.X}\%
  & \placeholder{XXX,XXX} & \placeholder{XXX,XXX}
  & \placeholder{XX,XXX} & \placeholder{XX,XXX} \\
L3 & \placeholder{X,XXX,XXX} & (ref.)
  & \placeholder{XXX,XXX} & \placeholder{XXX,XXX}
  & \placeholder{XX,XXX} & \placeholder{XX,XXX} \\
\bottomrule
\end{tabular}
\end{table*}

The L1$\rightarrow$L2 gap arises from a single mechanism: L1 does not
require generation to compensate network losses. At the Stadtbach
network's loss fraction of approximately \placeholder{X}\%, L1
systematically underestimates total heat generation requirements and
thus fuel consumption, electricity demand for heat pumps, and peak
grid import for demand charge purposes. The L2$\rightarrow$L3 gap
arises from the resolution of spatially heterogeneous dispatch
constraints that L2's five-zone aggregation cannot represent.

\subsection{\ce{CO2} Emissions Comparison}
\label{sec:results_co2}

\begin{table}[htbp]
\centering
\caption{Annual \ce{CO2} emissions comparison and primary mechanism.}
\label{tab:co2}
\begin{tabular}{@{}lllp{5cm}@{}}
\toprule
Level & Total [t\ce{CO2}/a] & $\Delta$ vs.\ L3 & Primary mechanism \\
\midrule
L1 & \placeholder{X,XXX} & $-$\placeholder{X.X}\%
  & HP over-dispatch; missing loss-driven boiler activation \\
L2 & \placeholder{X,XXX} & $-$\placeholder{X.X}\%
  & Aggregated spatial constraints \\
L3 & \placeholder{X,XXX} & (ref.) & -- \\
\bottomrule
\end{tabular}
\end{table}

\subsection{Dispatch Pattern Analysis}
\label{sec:results_dispatch}

\placeholder{[Figure~1: Load duration curves by generation source per
topology level. Three panels: heat pump output, gas boiler output,
electrode boiler output.]}

Key annual dispatch metrics:
\begin{itemize}
    \item Heat pump full-load hours: L1 overestimates by
    \placeholder{X}\% relative to L3.
    \item Gas boiler activation hours: L1 underestimates by
    \placeholder{X}\%; loss-driven peak events requiring boiler
    support are absent in L1.
    \item Electrode boiler utilization: relatively stable across
    levels (primarily price-driven).
\end{itemize}

\subsection{Storage Dispatch: The Primary Diagnostic Variable}
\label{sec:results_storage}

\begin{table}[htbp]
\centering
\caption{Storage dispatch metrics across topology levels. Charging
hour shift is the difference in mean charging start time relative
to L3 (winter, October--March).}
\label{tab:storage}
\begin{tabular}{@{}lllll@{}}
\toprule
Metric & L1 & L2 & L3 & Unit \\
\midrule
Annual equivalent full cycles   & \placeholder{XXX}   & \placeholder{XXX}   & \placeholder{XXX}   & -- \\
Mean charging hour (winter)     & \placeholder{XX:00} & \placeholder{XX:00} & \placeholder{XX:00} & h \\
Charging hour shift vs.\ L3    & \placeholder{+X.X}  & \placeholder{+X.X}  & 0                   & h \\
Peak shaving effectiveness      & \placeholder{XX}\%  & \placeholder{XX}\%  & \placeholder{XX}\%  & -- \\
Mean SOC at 06:00 (winter)     & \placeholder{XX}\%  & \placeholder{XX}\%  & \placeholder{XX}\%  & \% of $\bar{E}_s$ \\
\bottomrule
\end{tabular}
\end{table}

\placeholder{[Figure~2: Mean winter-week SOC profile, L1 vs.\ L3.
Key feature: L1 charges later (overnight minimum price); L3 charges
earlier (anticipating loss-amplified afternoon discharge cost).]}

\subsection{Network Losses and Computational Performance}
\label{sec:results_losses}

\begin{table}[htbp]
\centering
\caption{Annual thermal losses and computational performance.}
\label{tab:losses_computation}
\begin{tabular}{@{}lllll@{}}
\toprule
Level & Annual loss [MWh] & Loss fraction [\%] & Variables & Solve time [s] \\
\midrule
L1 & 0 (by design)        & 0                    & \placeholder{X,XXX}  & \placeholder{XX} \\
L2 & \placeholder{X,XXX}  & \placeholder{X.X}    & \placeholder{XX,XXX} & \placeholder{XXX} \\
L3 & \placeholder{X,XXX}  & \placeholder{X.X}    & \placeholder{XX,XXX} & \placeholder{X,XXX} \\
\bottomrule
\end{tabular}
\end{table}

%% ============================================================
%% 6. GENERALIZABILITY
%% ============================================================
\section{Generalizability: Synthetic Network Analysis}
\label{sec:generalizability}

\subsection{L1$\rightarrow$L2 Gap vs.\ Network Properties}
\label{sec:gen_L1L2}

The physics-inclusion gap is driven primarily by total pipe length and
associated annual loss energy.

\placeholder{[Figure~3a: Scatter/heatmap: L1$\rightarrow$L2 cost gap
[\%] vs.\ total pipe length [km], colored by storage ratio.
Expected: near-linear positive correlation.]}

For pipe lengths below 1~km, annual losses are $<$1\% of demand and
the L1$\rightarrow$L2 gap is economically negligible ($<$0.5\%). For
pipe lengths of 5--15~km, the gap grows to \placeholder{X--Y}\%.

\subsection{L2$\rightarrow$L3 Gap vs.\ Demand Heterogeneity}
\label{sec:gen_L2L3}

The spatial-refinement gap is driven by demand heterogeneity HI.
For HI~$<$~0.3, the gap remains below \placeholder{1}\% across all
synthetic configurations. For HI~$>$~0.5, the gap grows substantially,
reaching \placeholder{X--Y}\% for networks with storage ratio~$>$~6~h.

\placeholder{[Figure~3b: Heatmap: L2$\rightarrow$L3 gap [\%] as
function of HI (x-axis) and storage ratio (y-axis), medium pipe
length.]}

The interaction between heterogeneity and storage ratio is notable:
large storage magnifies spatial dispatch errors because incorrect
charging patterns compound over multiple daily cycles.

\subsection{Actionable Topology Selection Criteria}
\label{sec:gen_criteria}

\begin{table}[htbp]
\centering
\caption{Recommended topology selection criteria. Cost threshold for
``warranted'' is defined as topology gap $>$\placeholder{X}\% of L3
annual cost.}
\label{tab:criteria}
\begin{tabular}{@{}p{1.3cm}p{5.5cm}p{4.5cm}@{}}
\toprule
Level & Network conditions & Rationale \\
\midrule
\textbf{L1} & Total pipe length $<$1~km \textbf{AND} no thermal storage
  & Annual losses $<$1\% of demand; storage temporal arbitrage absent \\[4pt]
\textbf{L2} & Pipe 1--10~km \textbf{OR} storage ratio $\leq$6~h
  \textbf{AND} HI~$\leq$~0.4
  & Losses significant; spatial distribution approximately uniform \\[4pt]
\textbf{L3} & Pipe $>$10~km \textbf{OR} HI~$>$~0.5 \textbf{OR}
  storage ratio $>$6~h with HI~$>$~0.3
  & Spatial dispatch differences exceed cost threshold \\
\bottomrule
\end{tabular}
\end{table}

%% ============================================================
%% 7. SENSITIVITY ANALYSIS
%% ============================================================
\section{Sensitivity Analysis}
\label{sec:sensitivity}

\subsection{Parameter Sensitivity of L3}
\label{sec:sens_tornado}

\placeholder{[Figure~4: Tornado diagram for L3 annual operational cost.
Parameters: gas price $\pm$20\%, electricity price $\pm$20\%,
CO$_2$ price $\pm$50\%, COP $\pm$10\%, storage self-discharge $\pm$50\%,
ground temperature $\pm$3\,K.]}

\subsection{Stability of the Topology Gap Across Scenarios}
\label{sec:sens_gap}

\begin{table}[htbp]
\centering
\caption{Stability of topology gaps across economic and climate
scenarios (Stadtbach case study, L3 as reference).}
\label{tab:gap_stability}
\begin{tabular}{@{}lllll@{}}
\toprule
Scenario & L1--L3 gap & L2--L3 gap & L1 sign stable? & L2 sign stable? \\
\midrule
Baseline                & \placeholder{X.X\%} & \placeholder{X.X\%} & -- & -- \\
High gas                & \placeholder{X.X\%} & \placeholder{X.X\%} & \placeholder{Yes} & \placeholder{Yes} \\
Low electricity         & \placeholder{X.X\%} & \placeholder{X.X\%} & \placeholder{Yes} & \placeholder{Yes} \\
High \ce{CO2} price     & \placeholder{X.X\%} & \placeholder{X.X\%} & \placeholder{Yes} & \placeholder{Yes} \\
Cold year (+3\,K)       & \placeholder{X.X\%} & \placeholder{X.X\%} & \placeholder{Yes} & \placeholder{Yes} \\
Warm year ($-$3\,K)     & \placeholder{X.X\%} & \placeholder{X.X\%} & \placeholder{Yes} & \placeholder{Yes} \\
\bottomrule
\end{tabular}
\end{table}

%% ============================================================
%% 8. DISCUSSION
%% ============================================================
\section{Discussion}
\label{sec:discussion}

\subsection{When Does Topology Detail Matter?}
\label{sec:disc_when}

The dominant finding is that the L1$\rightarrow$L2 transition (physics
inclusion) drives the majority of the topology effect. For systems with
network losses constituting $>$3--5\% of annual heat supply---typical
for Central European DH networks with pipe lengths above 3~km---omitting
losses in L1 produces systematic cost underestimation of
\placeholder{X--Y}\%. This is not a modeling artefact but a direct
physical consequence: L1 requires less generation to meet the same
demand.

The L2$\rightarrow$L3 transition (spatial refinement) produces
substantially smaller differences for networks with uniform demand
(HI~$<$~0.4). If all demand nodes have similar load fractions, zonal
aggregation produces similar marginal generation cost signals, and
dispatch decisions at L2 approximate those at L3 closely. For the large
class of systems with moderate pipe length and near-uniform demand, L2
provides a favorable trade-off between fidelity and tractability. L1
should be reserved for networks where losses are demonstrably negligible.

\subsection{Storage as Diagnostic Variable and Planning Tool}
\label{sec:disc_storage}

The charging hour shift---L1 charges \placeholder{X}~hours later than
L3 on average in winter---constitutes a directly observable diagnostic
for topology adequacy in systems with thermal storage. A practitioner
using a copperplate model for operational planning would systematically
schedule late overnight charging. When the physical network incurs
higher losses during afternoon discharge, the true optimal operating
point requires earlier pre-charging. The economic cost of this dispatch
error is \placeholder{€X,XXX--XX,XXX}/year---a recurring annual loss
avoidable with an L2 model at negligible additional computational overhead
(factor \placeholder{X} increase in solve time).

\subsection{\ce{CO2} Policy Implications}
\label{sec:disc_co2}

Simplified topology models systematically underestimate \ce{CO2}
emissions for two compounding reasons. First, they over-dispatch the
heat pump (low \ce{CO2} source) because the additional generation for
loss compensation is invisible. Second, they miss boiler activation
events at peak-loss hours. For industrial networks subject to carbon
reporting obligations~\cite{EUheating2016}, L1-based emission accounting
introduces a systematic underestimation of
\placeholder{X--Y}~t\ce{CO2}/year.

\subsection{Relation to Existing Frameworks}
\label{sec:disc_frameworks}

The present findings indicate that existing energy system frameworks
using copperplate heat sector representation (oemof, PyPSA, FINE,
Calliope) will systematically underestimate operational costs and
emissions for DH systems with significant pipe length. The magnitude
depends on the loss fraction. For national-scale studies where heating
is one of many sectors, the additional complexity of an L2 representation
may not be warranted. For district-scale planning studies with explicit
component dispatch, L2 is the recommended minimum.

\subsection{Limitations}
\label{sec:discussion_limitations}

\begin{itemize}
    \item \textbf{Steady-state thermal model}: No transient startup
    dynamics or temperature ramping. Relevant for systems with frequent
    unit cycling.
    \item \textbf{No hydraulic pressure coupling}: Pumping energy not
    optimized; estimated at $<$2\% of network thermal losses at peak flow.
    \item \textbf{Hourly temporal resolution}: Sub-hourly demand peaks
    and price spikes not captured.
    \item \textbf{Single real-world case}: Generalization supported by
    the synthetic parametric analysis but not by multiple validation cases.
    \item \textbf{Deterministic framework}: Uncertainty addressed through
    sensitivity analysis, not stochastic or robust optimization.
    \item \textbf{Tree topology only}: Radial network structure assumed.
    Ring networks require additional flow direction binary variables and
    are not addressed here.
    \item \textbf{Dispatch-only horizon}: Potential topology effects on
    capacity sizing are not examined.
\end{itemize}

%% ============================================================
%% 9. CONCLUSION
%% ============================================================
\section{Conclusion}
\label{sec:conclusion}

This paper presents a controlled experimental comparison isolating
topology abstraction as the sole variable in district heating MILP
dispatch optimization. Three levels (L1 copperplate, L2 simplified
multi-node, L3 detailed multi-node) are implemented within a unified
framework with identical physics at L2 and L3, validated against an
Austrian industrial heat network.

\textbf{Principal findings:}
\begin{enumerate}
    \item The L1$\rightarrow$L2 transition (physics inclusion) dominates
    the topology effect, driving operational cost underestimation of
    \placeholder{X}\% and \ce{CO2} underestimation of \placeholder{X}\%.
    The mechanism is loss-driven generation underestimation.

    \item The L2$\rightarrow$L3 transition (spatial refinement) is
    marginal ($<$\placeholder{X}\%) for networks with HI~$<$~0.4 but
    grows substantially for concentrated demand distributions
    (HI~$>$~0.5).

    \item Thermal storage charging hour distribution is the most
    sensitive dispatch indicator, shifted by \placeholder{X}~hours
    between L1 and L3 in winter due to loss-dependent temporal holding
    costs invisible to L1.

    \item The L1$\rightarrow$L2 gap scales with total pipe length
    ($R^2\!=\,$\placeholder{XX}); the L2$\rightarrow$L3 gap scales
    with demand heterogeneity. These relationships yield actionable
    selection criteria.
\end{enumerate}

\textbf{Recommendations for practitioners:}
\begin{itemize}
    \item Use L2 as the minimum fidelity for any system with thermal
    storage or pipe length $>$1~km.
    \item Use L3 when demand heterogeneity index $>$~0.4 or storage
    ratio~$>$~6~h with heterogeneous demand.
    \item Validate storage dispatch patterns as a topology adequacy
    diagnostic before reporting cost or emission results.
    \item Conduct sensitivity analysis using the same topology level
    as the final model.
\end{itemize}

\textbf{Future research priorities:} Extension to ring topology
networks; stochastic and robust optimization variants; investment--dispatch
co-optimization to assess topology effects on capacity sizing; transient
validation against sub-hourly measurements.

%% ============================================================
%% REFERENCES
%% ============================================================
\bibliographystyle{elsarticle-num}
\begin{thebibliography}{99}

\bibitem{EUheating2016}
European Commission, An EU Strategy on Heating and Cooling,
COM(2016) 51 final, Brussels, 2016.

\bibitem{Lund2014_4GDH}
H.\ Lund, S.\ Werner, R.\ Wiltshire, S.\ Svendsen, J.E.\ Thorsen,
F.\ Hvelplund, B.V.\ Mathiesen,
4th Generation District Heating (4GDH): Integrating smart thermal
grids into future sustainable energy systems,
Energy 68 (2014) 1--11.

\bibitem{Werner2017}
S.\ Werner,
District heating and cooling in Sweden,
Energy 126 (2017) 419--429.

\bibitem{Frederiksen2013}
S.\ Frederiksen, S.\ Werner,
District Heating and Cooling, Studentlitteratur, Lund, 2013.

\bibitem{Talebi2016}
B.\ Talebi, P.A.\ Mirzaei, A.\ Bastani, F.\ Haghighat,
A Review of District Heating Systems: Modeling and Optimization,
Frontiers in Built Environment 2 (2016) 22.

\bibitem{Buffa2019}
S.\ Buffa, M.\ Canova, M.\ D'Antoni, M.\ Baratieri, R.\ Fedrizzi,
5th generation district heating and cooling systems: A review of
existing cases in Europe,
Renewable and Sustainable Energy Reviews 104 (2019) 504--522.

\bibitem{Lund2017_SmartEnergy}
H.\ Lund, P.A.\ Østergaard, D.\ Connolly, B.V.\ Mathiesen,
Smart energy and smart energy systems,
Energy 137 (2017) 556--565.

\bibitem{Fang2016}
T.\ Fang, R.\ Lahdelma,
Optimization of combined heat and power production with heat storage
based on sliding time window method,
Applied Energy 162 (2016) 723--732.

\bibitem{Guelpa2019review}
E.\ Guelpa, V.\ Verda,
Thermal energy storage in district heating and cooling systems:
A review,
Applied Energy 252 (2019) 113474.

\bibitem{Hilpert2018_oemof}
S.\ Hilpert, C.\ Kaldemeyer, U.\ Krien, S.\ Günther,
C.\ Wingenbach, G.\ Plessmann,
The Open Energy Modelling Framework (oemof) -- A new approach to
facilitate open science in energy system modelling,
Energy Strategy Reviews 22 (2018) 16--25.

\bibitem{Krien2020_oemofthermal}
U.\ Krien, S.\ Günther, J.\ Launer, D.\ Möller,
oemof.thermal: thermal components for the Open Energy Modelling
Framework, GitHub repository, 2020.

\bibitem{Brown2018_pypsa}
T.\ Brown, J.\ Hörsch, D.\ Schlachtberger,
PyPSA: Python for Power System Analysis,
Journal of Open Research Software 6 (2018) 4.

\bibitem{Welder2018_FINE}
L.\ Welder, D.\ Ryberg, L.\ Kotzur, T.\ Grube, M.\ Robinius,
D.\ Stolten,
Spatio-temporal optimization of a future energy system for
power-to-hydrogen applications in Germany,
Energy 158 (2018) 1130--1149.

\bibitem{Pfenninger2018_calliope}
S.\ Pfenninger, B.\ Pickering,
Calliope: a multi-scale energy systems modelling framework,
Journal of Open Source Software 3 (2018) 825.

\bibitem{Dorfner2016_urbs}
J.\ Dorfner,
urbs: A linear optimisation model for distributed energy systems,
GitHub repository, 2016.

\bibitem{Wiese2018_balmorel}
F.\ Wiese et al.,
Balmorel open source energy system model,
Energy Strategy Reviews 20 (2018) 26--34.

\bibitem{Connolly2014_EnergyPLAN}
D.\ Connolly, H.\ Lund, B.V.\ Mathiesen,
Smart Energy Europe: The technical and economic impact of one
potential 100\% renewable energy scenario for the European Union,
Renewable and Sustainable Energy Reviews 60 (2016) 1634--1653.

\bibitem{Vesterlund2017}
M.\ Vesterlund, A.\ Toffolo, J.\ Dahl,
Optimization of multi-source complex district heating network,
the case of Kiruna, Sweden,
Energy 126 (2017) 53--63.

\bibitem{Gabrielli2018}
P.\ Gabrielli, M.\ Gazzani, E.\ Martelli, M.\ Mazzotti,
Optimal design of multi-energy systems with seasonal storage,
Applied Energy 219 (2018) 408--424.

\bibitem{Wirtz2020}
M.\ Wirtz, L.\ Kivilip, P.\ Remmen, D.\ Müller,
5th Generation District Heating: A novel design approach based on
non-linear optimization,
Applied Energy 260 (2020) 114158.

\bibitem{Mertz2016}
T.\ Mertz, S.\ Serra, A.\ Henon, J.C.\ Reneaume,
A MINLP optimization of the configuration and the design of a
district heating network: Academic and real study cases,
Energy 116 (2016) 1156--1172.

\bibitem{Buenning2020}
F.\ Bünning, M.\ Wetter, M.\ Fuchs, D.\ Müller,
Bidirectional low temperature district energy systems with agent-based
control: Performance comparison and operation optimization,
Applied Energy 209 (2018) 502--515.

\bibitem{Giraud2015}
L.\ Giraud, M.\ Merabet, R.\ Bavière, M.\ Vallée,
Optimal control of district heating systems using dynamic simulation
and mixed integer linear programming,
Proceedings of the 12th International Modelica Conference, 2015.

\bibitem{Schweiger2017}
G.\ Schweiger, P.-O.\ Larsson, F.\ Magnusson, P.\ Lauenburg,
S.\ Velut,
District heating and cooling systems -- Framework for Modelica-based
simulation and dynamic optimization,
Energy 137 (2017) 566--578.

\bibitem{vanderHeijde2017}
B.\ van der Heijde, M.\ Fuchs, C.\ Ribas Tugores, G.\ Schweiger,
K.\ Sartor, D.\ Basciotti, D.\ Müller, C.\ Nytsch-Geusen,
M.\ Wetter, L.\ Helsen,
Dynamic equation-based thermo-hydraulic pipe model for district
heating and cooling systems,
Energy Conversion and Management 151 (2017) 158--169.

\bibitem{Blommaert2020}
M.\ Blommaert, J.\ Wack, M.\ Baelmans,
An adjoint optimization approach for the topological design of
large-scale district heating networks based on nonlinear models,
Applied Energy 280 (2020) 116025.

\bibitem{Paalsson1999}
H.\ Pálsson, A.\ Larsen, B.N.\ Jørgensen, S.\ Bøhm, Y.\ Zhou,
Equivalent models of district heating systems for on-line
optimization -- determination of model coefficients,
Proceedings of the 7th International Symposium on District Heating
and Cooling, Lund, 1999.

\bibitem{Duquette2016}
J.\ Duquette, A.\ Rowe, P.\ Wild,
Thermal performance of a steady state physical pipe model for
simulating district heating grids with variable flow,
Applied Energy 178 (2016) 383--393.

\bibitem{Vandermeulen2018}
A.\ Vandermeulen, B.\ van der Heijde, L.\ Helsen,
Controlling district heating and cooling networks to unlock
flexibility: A review,
Energy 151 (2018) 103--115.

\bibitem{Christidis2017}
A.\ Christidis, C.\ Koch, L.\ Pottel, G.\ Tsatsaronis,
The contribution of heat storage to the profitable operation of
combined heat and power plants in liberalized electricity markets,
Energy 41 (2012) 75--82.

\bibitem{Sartor2017}
K.\ Sartor, P.\ Dewallef,
Optimization of heat supply from combined heat and power plants
integrated in a district heating network,
Energy 124 (2017) 503--515.

\bibitem{Noussan2017}
M.\ Noussan, M.\ Jarre, A.\ Poggio,
Real operation data analysis on district heating load patterns,
Energy 129 (2017) 70--78.

\bibitem{pyomo}
W.E.\ Hart, C.\ Laird, J.-P.\ Watson, D.L.\ Woodruff,
Pyomo -- Optimization Modeling in Python, 2nd ed.,
Springer, 2017.

\bibitem{Paper2_CALION}
[Authors], CALION: A Flexible MILP Framework for Multi-Technology
District Heating Networks with Rolling Horizon Optimization,
[companion paper, in preparation, 2024].

\bibitem{Zenodo_data}
[Authors], Replication dataset: Topology abstraction effects on
district heating MILP dispatch, Zenodo [dataset], 2024.
\url{https://doi.org/10.5281/zenodo.XXXXXXX}

\end{thebibliography}

%% ============================================================
%% APPENDICES
%% ============================================================
\appendix

\section{Electricity Selling Price Formula}
\label{app:selling_price}

The complete electricity selling price, reflecting bilateral trading
conditions applicable to the Austrian industrial partner:
\begin{equation}
\lambda_t^{\mathrm{sell}} = \max\!\Bigg(\!
\underbrace{
\max\!\big(\lambda_t^{\mathrm{spot}} - \lambda^{\mathrm{spread}},\;
\lambda^{\mathrm{floor}}\big)
}_{\text{spread/floor-limited price}}
\cdot (1 - h^{\mathrm{sell}})
- \lambda^{\mathrm{sell,fee}} + \lambda^{\mathrm{sell,prem}},
\;\; 0\Bigg)
\end{equation}

Parameters: $\lambda^{\mathrm{spread}}$ is the bid-ask spread [\euro/MWh];
$\lambda^{\mathrm{floor}}$ a contractual minimum price floor;
$h^{\mathrm{sell}} \in [0,1]$ the selling haircut fraction representing
transaction costs and imbalance risk; $\lambda^{\mathrm{sell,fee}}$ the
export grid fee; $\lambda^{\mathrm{sell,prem}}$ a balancing service
premium, if applicable. The outer $\max(\cdot, 0)$ ensures the model
never incurs a financial penalty for selling.

\section{COP Bilinear Interpolation Error}
\label{app:cop}

The heat pump COP is pre-computed from a 2D manufacturer performance
map using bilinear interpolation on a grid with spacing
$\Delta T_s = \Delta T_a = 5$~K. For bilinear interpolation, errors
along the two axes are independent. The interpolation error is bounded
by the sum of the two univariate second-order terms:
\begin{equation}
|\mathrm{COP} - \widehat{\mathrm{COP}}| \leq
\frac{(\Delta T_s)^2}{8}
\max \left|\frac{\partial^2 \mathrm{COP}}{\partial T_s^2}\right|
+ \frac{(\Delta T_a)^2}{8}
\max \left|\frac{\partial^2 \mathrm{COP}}{\partial T_a^2}\right|
\label{eq:cop_error}
\end{equation}

For the heat pump models used in the case study, maximum second
derivatives are estimated from the manufacturer data sheet, yielding
an absolute COP error of $\pm$0.08--0.12 ($\approx$2--3\% relative
at typical COP~$\approx$~3.5).

Monte Carlo error propagation (1000 samples, COP perturbed independently
within error bounds at each timestep): system cost variation has a 95\%
confidence interval of $\pm$1.2\%. This is substantially smaller than
the L1--L3 topology gap, confirming that COP interpolation uncertainty
does not confound the topology comparison.

COP clamping to $[0.5, 15.0]$ activates in $<$0.3\% of timesteps
(extreme temperature conditions at system boundaries) and has negligible
effect on annual cost ($<$0.05\%).

\section{Anonymized Stadtbach Network Parameters}
\label{app:stadtbach}

Specific numerical parameters have been anonymized per agreement
with the industrial partner. \Cref{tab:app_params} provides the
aggregate network characteristics sufficient for reproducibility
verification.

\begin{table}[htbp]
\centering
\caption{Anonymized aggregate parameters, Stadtbach case study.}
\label{tab:app_params}
\begin{tabular}{@{}lll@{}}
\toprule
Parameter & Value & Unit \\
\midrule
Demand nodes (L3)                      & 30                        & -- \\
Demand zones (L2)                      & 5                         & -- \\
Total supply+return pipe length        & \placeholder{X.X}         & km \\
Pipe type distribution (DN100/200/250) & \placeholder{XX/XX/XX}    & \% of total length \\
Mean insulation $U$-value (supply)     & \placeholder{X.XX}        & W/(m$\cdot$K) \\
Mean insulation $U$-value (return)     & \placeholder{X.XX}        & W/(m$\cdot$K) \\
Design supply temperature              & \placeholder{XX}          & \si{\celsius} \\
Design return temperature              & \placeholder{XX}          & \si{\celsius} \\
Ground temperature (annual mean)       & \placeholder{X.X}         & \si{\celsius} \\
Ground temperature (seasonal range)   & \placeholder{X.X--X.X}    & \si{\celsius} \\
Peak aggregate demand                  & \placeholder{XX}          & MW \\
Annual aggregate demand                & \placeholder{XX}          & GWh \\
Demand heterogeneity index HI          & \placeholder{X.XX}        & -- \\
Storage energy capacity $\bar{E}_s$    & \placeholder{XXX}         & MWh \\
Storage charging ratio                 & \placeholder{X.X}         & h of peak \\
HP installed capacity                  & \placeholder{X.X}         & MW\textsubscript{th} \\
Gas boiler installed capacity          & \placeholder{XX}          & MW\textsubscript{th} \\
Electrode boiler installed capacity    & \placeholder{X.X}         & MW\textsubscript{th} \\
Grid import connection capacity        & \placeholder{XX}          & MW \\
Grid export connection capacity        & \placeholder{XX}          & MW \\
\bottomrule
\end{tabular}
\end{table}

\end{document}